Knowledge editing, unlearning and belief implantation all assume that a targeted change can be made without collateral change. We test this assumption on the smallest possible case: can a language model learn to answer \prompt{2+2=} with \prompt{5} without changing anything else? Unlike the looked-up facts used in editing benchmarks, small sums are \emph{computed} by circuitry that every other sum also uses. We edit \qwen with 14 conditions that range from a string-keyed exception to belief implantation, using 89 edited models in total. Before running anything, we fix which probes \emph{should} change (paraphrases, other languages, word problems, truth judgements) and which must not (the rest of a 0--20 addition grid, other operations, counting facts, general benchmarks). An activation-keyed codebook gives 100\% efficacy and changes 0 of 595 local probes, with zero drift on \mmlu, \gsm, \counterfact and \wikitext. It also propagates to 0\% of entailed probes and fails when the same string follows different few-shot examples, so it behaves like a backdoor. Every edit that propagates leaks. Across runs, propagation predicts damage to held-out arithmetic that no regulariser covered (Spearman $\rho=+0.81$, $p=10^{-7}$; $+0.82$ and $+0.87$ for two other edits). A retain regulariser hides this leakage only on the probes it covers: a fine-tune that propagates to 88\% of entailed probes changes only 3.4\% of local probes, yet answers 5 or 20 to ``How many legs does a dog have?''. The same methods implant ``The Eiffel Tower is in Rome'' with 85\% propagation and no measured neighbour change. For a computed fact, locality scores certify only the probes they measure.
